% !TeX root = Lec06_MIU.tex

\ifdefined\AdvMacroMain
  \chapter{Lecture 6: Money in Utility Models}
  \let\ThisNoteEnd\relax
\else
  \documentclass[12pt]{ctexart}
  \newcommand{\ProjectRoot}{..}
  \usepackage[a4paper,margin=0.85in]{geometry}
\usepackage{newpxtext}
\usepackage{amsmath,amssymb,amsthm,mathtools}
\usepackage{newpxmath}
\usepackage{xcolor}
\usepackage{enumitem}
\usepackage{graphicx}
\usepackage{array,tabularx}
\usepackage{float}
\usepackage{hyperref}

\definecolor{noteRed}{RGB}{190,45,90}

\newcommand{\red}[1]{\textcolor{noteRed}{#1}}
\newcommand{\blue}[1]{\textcolor{blue}{#1}}
\newcommand{\purple}[1]{\textcolor{purple}{#1}}

\newcommand{\E}{\mathbb{E}}
\newcommand{\dd}{\mathrm{d}}
\newcommand{\MPN}{\mathrm{MPN}}
\newcommand{\MRS}{\mathrm{MRS}}
\newcommand{\MRT}{\mathrm{MRT}}
\newcommand{\Var}{\operatorname{Var}}
\newcommand{\boxit}[1]{\fbox{\ensuremath{#1}}}
\newcommand{\circled}[1]{\textcircled{\scriptsize #1}}

\newenvironment{tightalign}
  {\[\begin{aligned}}
  {\end{aligned}\]}

\graphicspath{
  {\ProjectRoot/_assets/figures/}
  {\ProjectRoot/_assets/figures/lecture_04/}
}

  \title{Lecture 6: Money in Utility Models}
  \author{Chengyuan He}
  \date{April 2026}
  \begin{document}
  \maketitle
  \def\ThisNoteEnd{\end{document}}
\fi

\subsection*{Model I. Classical monetary model}
\begin{enumerate}[label=\arabic*]
\item assume perfect competition \& fully flexible prices
\item limited role for money,real and nominal dichotomy:
\[
N,C,Y,K,\frac{R}{P},\frac{M}{P}
\]
pinned down without monetary policy. Monetary policy only affects nominal variables.
\end{enumerate}

\paragraph{Household's problem}
\[
\max_{\{C_t,N_t,B_t\}}
E_0\sum_{t=0}^{\infty}\beta^t U(C_t,N_t)
\]

s.t.
\[
P_tC_t + Q_tB_t = B_{t-1}+W_tN_t+T_t
\]

where
\[
\begin{aligned}
B_t &: \text{one-period, nominally riskless discount bonds that mature in one period;}\\
&\quad \text{bought in } t,\ \text{and pay one unit of money at } t+1.\\[6pt]
Q_t &: \text{price of bonds; pay } Q_t \text{ units today and get 1 unit next period,} \\
&\quad \text{or, pay 1 unit today and get } \frac{1}{Q_t} \text{ units,i.e.gross return next period}.\\[6pt]
T_t &: \text{lump-sum taxes (if } T_t<0\text{) or dividends (if } T_t>0\text{).}
\end{aligned}
\]

Lagrangian:
\[
\mathcal L
=
E_0\sum_{t=0}^{\infty}\beta^t
\Big[
U(C_t,N_t)+\lambda_t(B_{t-1}+W_tN_t-T_t-P_tC_t-Q_tB_t)
\Big].
\]

FOCs:
\[
\frac{\partial \mathcal L}{\partial C_t}=0
\quad\Rightarrow\quad
U_{C_t}-\lambda_t P_t=0
\tag{1}
\]
\[
\frac{\partial \mathcal L}{\partial N_t}=0
\quad\Rightarrow\quad
U_{N_t}+\lambda_t W_t=0
\tag{2}
\]
\[
\frac{\partial \mathcal L}{\partial B_t}=0
\quad\Rightarrow\quad
-\beta^t\lambda_t Q_t+\beta^{t+1}E_t[\lambda_{t+1}]=0
\quad\Rightarrow\quad
Q_t=\beta E_t\left[\frac{\lambda_{t+1}}{\lambda_t}\right]
\tag{3}
\]
\[
\frac{\partial \mathcal L}{\partial \lambda_t}=0
\quad\Rightarrow\quad
P_tC_t+Q_tB_t=B_{t-1}+W_tN_t+T_t.
\tag{5}
\]

From the first two FOCs:
\[
-\frac{U_{N_t}}{U_{C_t}}
=
\frac{W_t}{P_t}.
\]

From the bond FOC:
\[
Q_t
=
\beta E_t\left[
\frac{U_{C,t+1}}{U_{C,t}}\frac{P_t}{P_{t+1}}
\right].
\]

Suppose
\[
U(C_t,N_t)
=
\frac{C_t^{1-\sigma}}{1-\sigma}
-
\frac{N_t^{1+\varphi}}{1+\varphi}.
\]

Then
\[
U_{C_t}=C_t^{-\sigma},
\qquad
U_{N_t}=-N_t^{\varphi}.
\]

So
\[
\frac{W_t}{P_t}=C_t^\sigma N_t^\varphi.
\]

Hence
\[
N_t=\left(\frac{W_t/P_t}{C_t^\sigma}\right)^{1/\varphi}.
\]
\[
Q_t
=
\beta E_t\!\left[
\left(\frac{C_{t+1}}{C_t}\right)^{-\sigma}
\frac{P_t}{P_{t+1}}
\right].
\]
\[
\frac{W_t}{P_t}=C_t^\sigma N_t^\varphi
\quad\Rightarrow\quad
N_t=\left(\frac{W_t/P_t}{C_t^\sigma}\right)^{1/\varphi}
\quad\Rightarrow\quad
\log N_t
=
\frac{1}{\varphi}\log\!\left(\frac{W_t}{P_t}\right)
-
\frac{\sigma}{\varphi}\log C_t
\quad\Rightarrow\quad
\frac{\partial \log N_t}{\partial \log (W_t/P_t)}
=
\frac{1}{\varphi}
\]

\noindent
\textcolor{red}{\(\dfrac{1}{\varphi}\): Frisch elasticity (how labor supply responds to real wage), \(\varphi \uparrow \Rightarrow\) elasticity \(\downarrow\).}
\paragraph{Firm's problem}

\[
\max_{Y_t,\,N_t}\ \Pi_t = P_t Y_t - W_t N_t
\]

s.t.
\[
Y_t=A_tN_t^{1-\alpha}.
\]

FOC:
\[
\frac{W_t}{P_t}=(1-\alpha)A_tN_t^{-\alpha}
\qquad
\text{: both are real terms.}
\tag{4}
\]

\paragraph{Market clearing}
\[
\qquad \text{Goods market clearing:}
\qquad Y_t=C_t
\]
\[
\text{Labor market clearing: }\qquad N_t^{(d)} = N_t^{(s)}
(\text{implicitly holds})
\]
\[
\text{Bond market clearing:Hold by Walras Law}
\qquad B_t=0
\]

Then we have the system
\[
\left\{
\begin{aligned}
Y_t &= C_t, 
& \qquad & \text{goods market clearing} \\[4pt]
T_t &= \alpha P_t Y_t, 
& & \text{dividend} \\[4pt]
Y_t &= A_t N_t^{1-\alpha}, 
& & \text{production function} \\[4pt]
\frac{W_t}{P_t} &= C_t^{\sigma} N_t^{\varphi}, 
& & \text{labor-leisure condition} \\[6pt]
\frac{W_t}{P_t} &= (1-\alpha)A_t N_t^{-\alpha}, 
& & \text{real wage = MPL} \\[6pt]
P_t C_t + Q_t B_t &= B_{t-1} + W_t N_t + T_t, 
& & \text{budget constraint} \\[6pt]
Q_t &= \beta E_t \left[
\left(\frac{C_{t+1}}{C_t}\right)^{-\sigma}
\frac{P_t}{P_{t+1}}
\right], 
& & \text{bond Euler equation} \\[8pt]
\log A_t &= (1-\rho_a)\log \bar{A} + \rho_a \log A_{t-1} + \varepsilon_t,
& & \text{productivity process}
\end{aligned}
\right.
\]
\[
\text{No monetary policy involved.}
\text{Denote the real wage as }
W_t^r = \frac{W_t}{P_t}.
\]

\[
\Rightarrow
\qquad
\text{Real equilibrium(in units of goods):}
\qquad
\left\{
\begin{aligned}
W_t^r &= C_t^{\sigma} N_t^{\varphi} \qquad (1)\\[4pt]
Y_t &= A_t N_t^{1-\alpha} \qquad (2)\\[4pt]
W_t^r &= (1-\alpha)A_t N_t^{-\alpha} \qquad (3)\\[4pt]
Y_t &= C_t \qquad (4)
\end{aligned}
\right.
\]

\[
\text{Four unknowns: } W_t^r,\ C_t,\ N_t,\ Y_t.
\qquad
\text{Four equations.}
\]
\paragraph{At steady state}
Assume
\[
\bar A=1.
\]

From (2), (3), and \(Y=C\):
\[
(1-\alpha)\bar N^{-\alpha}
=
\bar C^\sigma \bar N^\varphi.
\]

Since
\[
\bar Y=\bar N^{1-\alpha},
\qquad
\bar C=\bar Y,
\]
we have
\[
(1-\alpha)\bar N^{-\alpha}
=
\left(\bar N^{1-\alpha}\right)^\sigma \bar N^\varphi
=
\bar N^{\sigma(1-\alpha)+\varphi}.
\]

Hence
\[
1-\alpha
=
\bar N^{\sigma(1-\alpha)+\varphi+\alpha},
\]
so
\[
\bar N
=
(1-\alpha)^{\frac{1}{\sigma(1-\alpha)+\varphi+\alpha}}.
\]

Thus
\[
\bar Y=\bar N^{1-\alpha}
=
(1-\alpha)^{\frac{1-\alpha}{\sigma(1-\alpha)+\varphi+\alpha}},
\qquad
\bar C=\bar Y.
\]
and
\[
\bar{W}^{\,r}
=
\bar C^\sigma \bar N^\varphi
=
(1-\alpha)\bar N^{-\alpha}.
\]

\paragraph{Log-linearization}

\[
\hat w_t^r=\sigma \hat c_t+\varphi \hat n_t
\]
\[
\hat y_t=\hat a_t+(1-\alpha)\hat n_t
\]
\[
\hat w_t^r=\hat a_t-\alpha \hat n_t
\]
\[
\hat y_t=\hat c_t.
\]

Only one shock \(\hat a_t\), no predetermined variable. Assume
\[
\hat y_t=\psi_{ya}\hat a_t,
\qquad
\hat n_t=\psi_{na}\hat a_t,
\qquad
\hat{w}_t^{\,r}=\psi_{wa}\,\hat{a}_t
\]

Then
\[
\sigma \hat{a}_t + \sigma(1-\alpha)\hat{n}_t + \varphi \hat{n}_t
= \hat{a}_t - \alpha \hat{n}_t
\]

\[
\Rightarrow\quad
\hat{c}_t=\hat{y}_t
=
\frac{1+\varphi}{\sigma(1-\alpha)+\varphi+\alpha}\,\hat{a}_t
\qquad
\text{coefficient >0}
\]

\[
\hat{w}_t^{\,r}
=
\frac{\varphi+\sigma}{\sigma(1-\alpha)+\varphi+\alpha}\,\hat{a}_t
\qquad
\text{coefficient >0}
\]
\[
\hat{n}_t
=
\frac{1-\sigma}{\sigma(1-\alpha)+\varphi+\alpha}\,\hat{a}_t
\qquad
\text{coefficient is ambiguous, depends on } \sigma
\]
\textcolor{red}{Income effect dominates if \(\sigma>1\): \(\dfrac{\partial \hat n_t}{\partial \hat a_t}<0\).}

\textcolor{red}{Substitution effect dominates if \(\sigma<1\): \(\dfrac{\partial \hat n_t}{\partial \hat a_t}>0\).}

\paragraph{Real interest rate}

Bond Euler:
\[
Q_t
=
\beta E_t \left[
\left(\frac{C_{t+1}}{C_t}\right)^{-\sigma}
\frac{1}{\Pi_{t+1}}
\right].
\tag{*}
\]
\[
\left(\frac{C_{t+1}}{C_t}\right)^{-\sigma}
:\ \text{real stochastic discount factor},
\qquad
\frac{1}{\Pi_{t+1}}
:\ \text{converts } t+1 \text{ nominal payoff into real terms}.
\]

Let
\[
Q_t=e^{-i_t},
\qquad
\frac{P_{t+1}}{P_t}=\Pi_{t+1}\ \text{inflation},
\qquad
\pi_{t+1}
=
\log \Pi_{t+1}
=
\log P_{t+1}-\log P_t,
\qquad
\beta = e^{-\rho}
\]
\[
\text{Define real interest rate: }\quad
r_t = i_t - E_t\pi_{t+1}
\]
\[
i_t = -\log Q_t
\qquad \text{log nominal interest rate}
\]

\[
\rho = -\log \beta
\qquad \text{rate of time preference}
\]
\qquad
\[
\text{Define new deviation: }\quad
\hat{i}_t=i_t-\bar{i}
\qquad
\hat{r}_t=r_t-\bar{r}
\qquad
\bar{\pi}=\log \bar{P}-\log \bar{P}=0
\]
\qquad
\[
\left\{
\begin{aligned}
\left(\frac{C_{t+1}}{C_t}\right)^{-\sigma}
&=
e^{-\sigma \log C_{t+1}} \big/ e^{-\sigma \log C_t}
=
e^{-\sigma(\log C_{t+1}-\log C_t)},
\\[6pt]
\frac{1}{\Pi_{t+1}}
&=
\frac{P_t}{P_{t+1}}
=
e^{\log P_t-\log P_{t+1}}
=
e^{-\pi_{t+1}}.
\end{aligned}
\right.
\]

\[
\Rightarrow\quad
\textcolor{red}{
e^{-i_t}
=
e^{-\rho} E_t \!\left[
e^{-\sigma(\hat c_{t+1}-\hat c_t)-\pi_{t+1}}
\right]
}
\tag{**}
\]

\[
-i_t = -\rho + E_t\!\left[-\sigma(\hat c_{t+1}-\hat c_t)-\pi_{t+1}\right]
\tag{1}
\]

\[
\left.
\begin{aligned}
-\bar{i}
&=
-\rho
+
E_t\!\left[
-\sigma(\log \bar{C}-\log \bar{C})-\bar{\pi}
\right],
\\[6pt]
\bar{\pi}
&=
\log \bar{P}-\log \bar{P}
=
0
\end{aligned}
\right\}
\;\Rightarrow\;
\bar{i}=\rho
\tag{2}
\]

\[
\begin{aligned}
(1)+(2)\quad \Rightarrow\quad
-\hat{i}_t+\bar{i}
&=
E_t\!\left[-\sigma\bigl(\log C_{t+1}-\log C_t\bigr)-\pi_{t+1}\right].
\end{aligned}
\]

\[
\left\{
\begin{aligned}
\hat{c}_{t+1} &= \log C_{t+1}-\log \bar{C},\\
\hat{c}_t &= \log C_t-\log \bar{C}
\end{aligned}
\right.
\qquad \Rightarrow \qquad
\hat{c}_{t+1}-\hat{c}_t=\log C_{t+1}-\log C_t.
\]

\[
\begin{aligned}
\Rightarrow\quad
-\hat{i}_t
&=
E_t\!\left[-\sigma(\hat{c}_{t+1}-\hat{c}_t)-\pi_{t+1}\right],
\\[6pt]
\Rightarrow\quad
-\hat{i}_t
&=
-\sigma E_t\Delta \hat{c}_{t+1}-E_t[\pi_{t+1}],
\\[6pt]
\Rightarrow\quad
\sigma E_t\Delta \hat{c}_{t+1}
&=
\hat{i}_t-E_t[\pi_{t+1}],
\\[6pt]
\hat{\pi}_{t+1}
&=
\pi_{t+1}-\bar{\pi}
=
\pi_{t+1},
\\[6pt]
\Rightarrow\quad
\hat{i}_t-E_t[\hat{\pi}_{t+1}]
&=
\sigma E_t\Delta \hat{c}_{t+1}.
\end{aligned}
\tag{***}
\]
\[
\text{Since }
r_t=i_t-E_t[\pi_{t+1}],
\tag{3}
\]
\[
\bar{r}=\bar{i}-\bar{\pi}=\bar{i},
\tag{4}
\]
\[
(3)-(4)\;\Rightarrow\;
\hat{r}_t=\hat{i}_t-E_t[\hat{\pi}_{t+1}]
\]
\[
\text{Plug in }(***)\;\Rightarrow\;
\hat{r}_t=\sigma E_t\Delta \hat{c}_{t+1}
\]
\[
\hat{c}_{t+1}=\hat{y}_{t+1}=\psi_{ya}\hat{a}_{t+1},
\]
\[
\Rightarrow\;
\hat{r}_t=\sigma \psi_{ya}\, E_t[\Delta \hat{a}_{t+1}]
\]
\[
\text{If the shock is transitory / mean reverting}
\quad\Rightarrow\quad
E_t[\Delta \hat{a}_{t+1}]<0
\quad\Rightarrow\quad
r_t<0.
\]
\textcolor{red}{A positive transitory technology shock causes the current real interest rate to fall below its steady-state level.}

\textcolor{red}{Intuition: \(a_t\uparrow\) gives a favorable temporary income shock \(\Rightarrow\) households smooth consumption by saving \(\Rightarrow\) bond demand rises \(B_t^d\uparrow\) \(\Rightarrow\) bond price rises \(Q_t\uparrow\) \(\Rightarrow\) real interest rate falls \(r_t\downarrow\).}
\paragraph{Dichotomy}

\begin{itemize}
\item \(r_t\) pinned down by real side, money growth does not affect \(r_t\).
\item No policy is the best since only real variables enter utility.
\item Does not match empirics:data say real variable move after monetary shocks and nominall variables respond slugg.
\end{itemize}

\newpage
\subsection*{Model II. Money in utility model}

\begin{itemize}
\item Make nominal terms affect real terms.
\item Explain why agents hold money which is inferior to bond. (Money has no return but affects utility.)
\end{itemize}
\paragraph{Household's Problem}

\[
\max_{\{C_t,M_t,N_t\}_{t=0}^{\infty}}
E_0\sum_{t=0}^{\infty}\beta^t\,
U\!\left(C_t,\mathbb{M}_t,N_t\right)
\]

\[
\text{s.t.}\quad
P_t C_t + Q_t B_t + M_t \le B_{t-1} + M_{t-1} + W_t N_t + T_t
\]

\textcolor{red}{Let \(\mathbb{M}_t=\dfrac{M_t}{P_t}\), the real money balance/supply.}

Financial wealth before decisions(predetermined variable):
\[
A_t = B_{t-1}+M_{t-1}
\]

\[
\Lambda_{t,T}=Q_t\cdot Q_{t+1}\cdots Q_T
\]

B.C.:
\[
P_tC_t + Q_tA_{t+1} + (1-Q_t)M_t = A_t + W_tN_t + T_t
\]
Notes 1:
\[
i_t:\ \text{continuously compounded nominal rate}
\]
\[
Q_t:\ \text{price today } Q_t,\ \text{payoff tomorrow } 1
\]

Nominal interest rate:
\[
1-Q_t = 1-e^{-i_t} \approx i_t
\]
\[
\begin{aligned}
1-Q_t \;:\;& \text{the opportunity cost of holding one's financial wealth in money rather than} \\
& \text{in interest-bearing bonds; a high } i_t \text{ thus means contractionary monetary policy.}
\end{aligned}
\]
Notes:
\[
\Lambda_{t,T}:\ \text{state-contingent discount factor from } t \text{ to } T
\]
\[
\frac{A_T}{P_T}:\ \text{terminal real wealth}
\]

No-Ponzi scheme:
\[
\lim_{T\to\infty} E_t\!\left[\Lambda_{t,T}\frac{A_T}{P_T}\right] \ge 0,
\qquad \forall t
\]

Notes:
\[
\text{cannot end up with infinitely bad wealth}
\]
\[
\text{(borrow a lot} \Rightarrow \text{roll-over)}
\]

TVC:
\[
\lim_{T\to\infty} E_t\!\left[\Lambda_{t,T}\frac{A_T}{P_T}\right] = 0
\]

Note:
\[
\text{no resources on table when world ends}
\]

\[
\text{All financial assets } A_t \text{ yield gross nominal return } Q_t^{-1}=e^{i_t}.
\]



Lagrangian:
\[
\mathcal L
=
E_0\sum_{t=0}^{\infty}\beta^t
\left[
U\!\left(C_t,\mathbb{M}_t,N_t\right)
+
\lambda_t\big(
A_t+W_tN_t+T_t-P_tC_t-Q_tA_{t+1}-(1-Q_t)M_t
\big)
\right].
\]

FOCs:
\[
\frac{\partial \mathcal L}{\partial M_t}=0
\quad\Rightarrow\quad
U_{\mathbb{M}_t}\frac{1}{P_t}-\lambda_t(1-Q_t)=0
\]
\[
\frac{\partial \mathcal L}{\partial C_t}=0
\quad\Rightarrow\quad
U_{C,t}=\lambda_t P_t
\]
\[
\frac{\partial \mathcal{L}}{\partial N_t}=0
\quad\Rightarrow\quad
-\frac{U_{N,t}}{U_{C,t}}=\frac{W_t}{P_t}.
\]
Hence
\[
\frac{U_{\mathbb{M}_t}}{U_{C_t}}
=
1-Q_t
=
1-e^{-i_t}.
\tag{5}
\]


\subsection*{Separable utility case(still dichotomy)}

Assume
\[
U\!\left(C_t,\mathbb{M}_t,N_t\right)
=
\frac{C_t^{1-\sigma}}{1-\sigma}
+
\frac{(\mathbb{M}_t)^{1-\nu}}{1-\nu}
-
\frac{N_t^{1+\varphi}}{1+\varphi}.
\]

Neither \(U_{C_t}\) nor \(U_{N_t}\) depends on \(\mathbb{M}_t\)\(\Rightarrow\) real equilibrium not affected(same as before).

But we can derive money demand:
\[
U_{\mathbb{M}_t}
=
\left(\mathbb{M}_t\right)^{-\nu},
\qquad
U_{C_t}=C_t^{-\sigma}.
\]
\[
\frac{U_{\mathbb{M},t}}{U_{C,t}} = 1-e^{-i_t}
\quad\Rightarrow\quad
\frac{\mathbb{M}_t^{-\nu}}{C_t^{-\sigma}} = 1-e^{-i_t}
\]

\[
\Rightarrow\quad
\mathbb{M}_t
=
C_t^{\sigma/\nu}\left[1-e^{-i_t}\right]^{-1/\nu}
\]

\[
\Rightarrow\quad
\log M_t-\log P_t
=
\frac{\sigma}{\nu}\log C_t
-
\frac{1}{\nu}\log(1-e^{-i_t})
\]
\[
\hat{\mathbb{M}}_t
\approx
\frac{\sigma}{\nu}\hat{c}_t
-\frac{1}{\nu\left(e^{\bar{i}}-1\right)}\hat{i}_t.
\]
\text{math notes:}
\[
\begin{aligned}
\text{Let } f(i) &= \log(1-e^{-i}),\\[4pt]
f'(i) &= \frac{e^{-i}}{1-e^{-i}} = \frac{1}{e^i-1},\\[4pt]
f''(i) &= -\frac{e^i}{(e^i-1)^2},\\[6pt]
\log(1-e^{-i_t})
&\approx
\log(1-e^{-\bar{i}})
+\frac{1}{e^i-1}(i_t-\bar{i}).
\end{aligned}
\]
\[
\log(1-e^{-\bar{i}}) \text{ is constant, so it cancels out in derivation.}
\]
\mbox{}\\
If \(\sigma=\nu\), define
\[
\eta\equiv \frac1{\nu\left(e^{\bar{i}}-1\right)},
\]
then approximately
\[
\hat m_t-\hat p_t=\hat y_t-\eta \hat i_t.
\]

\textcolor{red}{\(\eta\): semi-elasticity of real money demand with respect to the nominal interest rate.}

\textcolor{red}{\(\hat m_t-\hat p_t\): real money balances are determined by \(\hat y_t\) and \(\hat i_t\).}

\textcolor{red}{\(\hat y_t\): determined by the real-side equilibrium, so monetary policy does not affect real activity in the separable case.}

\subsection*{Inseparable utility case}
Utility:
\[
U\!\left(C_t,\mathbb{M}_t,N_t\right)
=
\frac{X_t^{1-\sigma}}{1-\sigma}
-
\frac{N_t^{1+\varphi}}{1+\varphi}.
\]
\[
\text{where, }\;
X_t=
\begin{cases}
\left[(1-\theta)C_t^{\,1-\nu}+\theta\mathbb{M}_t^{1-\nu}\right]^{\frac{1}{1-\nu}}, & \nu\neq 1,\\[10pt]
C_t^{\,1-\theta}\mathbb{M}_t^{\theta}, & \nu=1.
\end{cases}
\]
\[
\text{Elasticity of substitution between } C_t \text{ and } \mathbb{M}_t
\text{ is } \frac{1}{\nu}
\]
\[
\theta\in(0,1)
\qquad
\text{weight on money inside the CES aggregator.}
\]

\text{Derivatives when \(\nu\neq 1\)}

For \(\nu\neq 1\),
\[
X_t=\left[(1-\theta)C_t^{\,1-\nu}+\theta \mathbb{M}_t^{\,1-\nu}\right]^{\frac{1}{1-\nu}}.
\]

Then
\[
U_{C,t}=X_t^{-\sigma}\cdot \frac{1}{1-\nu}\cdot X_t^{\nu}\cdot (1-\theta)(1-\nu)C_t^{-\nu}
\]
so
\[
U_{C,t}=(1-\theta)X_t^{\,\nu-\sigma}C_t^{-\nu}.
\]

Similarly,
\[
U_{\mathbb{M},t}
=
X_t^{-\sigma}\cdot \frac{1}{1-\nu}\cdot X_t^{\nu}\cdot \theta(1-\nu)\mathbb{M}_t^{-\nu}
\]
so
\[
U_{\mathbb{M},t}=\theta X_t^{\,\nu-\sigma}\mathbb{M}_t^{-\nu}.
\]

Also,
\[
U_{N,t}=-N_t^{\varphi}.
\]

\paragraph{Marginal rate of substitution and elasticity}

Define
\[
X_C=(1-\theta)X^\nu C^{-\nu},
\qquad
X_ \mathbb{M}=\theta X^\nu \mathbb{M}^{-\nu}.
\]

Then
\[
MRS_{C,\mathbb{M}}
=
\frac{U_C}{U_{\mathbb{M}}}
=
\frac{X_C}{X_{\mathbb{M}}}
=
\frac{1-\theta}{\theta}\left(\frac{\mathbb{M}}{C}\right)^\nu.
\]
Hence
\[
\varepsilon
=
\frac{d\ln(\mathbb{M}/C)}{d\ln MRS_{C,\mathbb{M}}}
=
\frac{1}{\nu}.
\]

Thus,
\[
\nu>1 \;\Rightarrow\; \varepsilon<1 \quad \text{(complements)},
\]
\[
\nu<1 \;\Rightarrow\; \varepsilon>1 \quad \text{(substitutes)}.
\]

\paragraph{FOCs}

Labor supply condition:
\[
-\frac{U_{N,t}}{U_{C,t}}=\frac{W_t}{P_t}.
\]

Substituting the derivatives,
\[
\left\{
\begin{aligned}
\frac{W_t}{P_t}
&=
\frac{1}{1-\theta}N_t^\varphi X_t^{\sigma-\nu}C_t^\nu, \\[8pt]
Q_t
&=
\beta E_t\!\left[\frac{U_{C,t+1}}{U_{C,t}}\frac{P_t}{P_{t+1}}\right]
=
\beta E_t\!\left[
\left(\frac{C_{t+1}}{C_t}\right)^{-\nu}
\left(\frac{X_{t+1}}{X_t}\right)^{\nu-\sigma}
\frac{P_t}{P_{t+1}}
\right].
\end{aligned}
\right.
\qquad
\begin{array}{l}
X_t \text{ depends on } \mathbb{M}_t,\\
\text{thus money is non-neutral.}
\end{array}
\]

Money FOC:
\[
\frac{U_{\mathbb{M},t}}{U_{C,t}}=1-e^{-i_t}.
\]

So
\[
\frac{\theta}{1-\theta}C_t^\nu \mathbb{M}_t^{-\nu}=1-e^{-i_t}.
\]

Solving for real money balances,
\[
\mathbb{M}_t
=
C_t\left(1-e^{-i_t}\right)^{-1/\nu}
\left(\frac{\theta}{1-\theta}\right)^{1/\nu}.
\]
When $\nu=\sigma$ and we drop $\theta$, the utility function reduces to the separable case, so the dichotomy holds.

When $\nu \neq \sigma$, we have the general (non-separable) case.
\[
\begin{aligned}
&\text{1) Labor supply and Euler equation are affected by the level of real balances through } X_t,\\
&\text{2) Real balance depends on nominal interest rate.}
\end{aligned}
\]

\[
\begin{aligned}
\text{Non-neutral monetary policy:}\qquad
&\text{Monetary policy}
\;\Rightarrow\;
\text{nominal rate } i_t\\
&\Rightarrow\;
\text{real balance } \mathbb{M}
\;\Rightarrow\;
\text{labor supply and output.}
\end{aligned}
\]
\[
\text{Log linearize:}\qquad
\mathbb{M}_t
=
C_t(1-e^{-i_t})^{-1/{\nu}}
\left(\frac{\theta}{1-\theta}\right)^{1/{\nu}}
\;\Rightarrow\;
\hat{m}_t-\hat{p}_t
=
\hat{c}_t-\eta \hat{i}_t
\qquad \cdots\; (*)
\]

\[
\frac{W_t}{P_t}
=
\frac{1}{1-\theta}
N_t^{\varphi}X_t^{\sigma-\nu}C_t^{\nu}
\;\Rightarrow\;
\hat{w}_t-\hat{p}_t
=
\sigma\hat{c}_t+\varphi \hat{n}_t+(\nu-\sigma)(\hat{c}_t-\hat{x}_t)
\qquad \cdots\; (**)
\]
\[
X_t^{1-\nu}
=
(1-\theta)C_t^{1-\nu}
+
\theta \mathbb{M}^{1-\nu}
\]
\[
\textcolor{red}{\text{S.S.:}\quad
(1-\theta)\bar{C}^{\,1-\nu}
=
\bar{X}^{\,1-\nu}
-
\theta \bar{\mathbb{M}}^{\,1-\nu}}
\]
\[
\Rightarrow\quad
\bar X^{1-\nu}\hat x_t
=
(1-\theta)\bar C^{1-\nu}\hat c_t
+
\theta \bar{\mathbb{M}}^{1-\nu}(\hat m_t-\hat p_t).
\]

\[
\text{Using } 
(1-\theta)\bar C^{1-\nu}
=
\bar X^{1-\nu}
-
\theta\bar{\mathbb{M}}^{1-\nu},
\]
\[
\bar X^{1-\nu}\hat x_t
=
\left[
\bar X^{1-\nu}
-
\theta\bar{\mathbb{M}}^{1-\nu}
\right]\hat c_t
+
\theta\bar{\mathbb{M}}^{1-\nu}(\hat m_t-\hat p_t).
\]

\[
\Rightarrow
\bar X^{1-\nu}(\hat c_t-\hat x_t)
=
\theta\bar{\mathbb{M}}^{1-\nu}
\left[
\hat c_t-(\hat m_t-\hat p_t)
\right].
\]

\[
\Rightarrow
\hat c_t-\hat x_t
=
\underbrace{
\theta
\left(
\frac{\bar{\mathbb{M}}}{\bar X}
\right)^{1-\nu}
}_{\chi}
\left[
\hat c_t-(\hat m_t-\hat p_t)
\right].
\tag{***}
\]


\[
\text{Define}\qquad
\left(\frac{\bar{M}}{\bar{P}}\right)\Big/\bar{C}
=
k_m
\]

\[
\textcolor{red}{k_m:\ \text{inverse consumption velocity, low }k_m\text{ high velocity}}
\]

\[
\text{Let}\qquad
\chi
=
\theta
\frac{\left(\frac{\bar{M}}{\bar{P}}\right)^{1-\nu}}{\bar{X}^{\,1-\nu}}
=
\theta
\frac{\left(\frac{\bar{M}}{\bar{P}}\right)^{1-\nu}}
{(1-\theta)\bar{C}^{\,1-\nu}+\theta\left(\frac{\bar{M}}{\bar{P}}\right)^{1-\nu}}
=
\frac{
\theta\left(\left(\frac{\bar{M}}{\bar{P}}\right)\big/\bar{C}\right)^{1-\nu}
}{
1-\theta+\theta\left(\left(\frac{\bar{M}}{\bar{P}}\right)\big/\bar{C}\right)^{1-\nu}
}
\]

\[
=
\frac{\theta k_m^{\,1-\nu}}{1-\theta+\theta k_m^{\,1-\nu}}
\]

\[
\textcolor{red}{\chi:\ \text{steady-state expenditure share of money inside the CES aggregator}}
\]

\[
\mathbb{M}_t
=
C_t(1-e^{-i_t})^{-1/{\nu}}
\left(\frac{\theta}{1-\theta}\right)^{1/{\nu}}
\;\Rightarrow\;
\left(\frac{\bar{M}}{\bar{P}}\right)\Big/\bar{C}
=
(1-\beta)^{-1/{\nu}}
\left(\frac{\theta}{1-\theta}\right)^{1/{\nu}}
\]

\[
\Rightarrow\quad
k_m
=
\left[
\frac{\theta}{(1-\beta)(1-\theta)}
\right]^{1/{\nu}}
\]

\[
\textcolor{red}{k_m \text{ increases with money weight } \theta,\ \beta\uparrow \Rightarrow \text{opportunity cost of holding money} \downarrow}
\]

\[
\text{Plug } (***) \text{ into } (**), \text{ we have}
\]

\[
\hat{w}_t-\hat{p}_t
=
\sigma\hat{c}_t+\varphi \hat{n}_t
+
(\nu-\sigma)\chi
\Bigl[\hat{c}_t-(\hat{m}_t-\hat{p}_t)\Bigr]
\]

\[
\text{Plug } (*) \text{ into the above equation:}
\]

\[
\hat{w}_t-\hat{p}_t
=
\sigma\hat{c}_t+\varphi \hat{n}_t
+
(\nu-\sigma)\chi
\Bigl[\hat{c}_t-\hat{y}_t+\eta \hat{i}_t\Bigr]
\]

\[
\Rightarrow\quad
\hat{w}_t-\hat{p}_t
=
\sigma\hat{c}_t+\varphi \hat{n}_t+\omega \hat{i}_t
\]

\[
\text{where}\qquad
\omega=(\nu-\sigma)\chi\eta
\]

\[
\textcolor{red}{\omega:\ \text{like nominal interest rate and labor supply}}
\]
\[
\chi
=
\frac{\theta^{1/\nu}(1-\beta)^{1-\frac{1}\nu}}
{(1-\theta)^{1/\nu}+\theta^{1/\nu}(1-\beta)^{1-\frac{1}\nu}}
=
\frac{\left(\frac{\theta}{1-\theta}\right)^{1/\nu}(1-\beta)^{1-\frac{1}\nu}}
{1+\left(\frac{\theta}{1-\theta}\right)^{1/\nu}(1-\beta)^{1-\frac{1}\nu}}
\]

\[
\Rightarrow\quad
\chi=\frac{k_m(1-\beta)}{1+k_m(1-\beta)}
\]

\[
k_m(1-\beta)
=
\left(\frac{\theta}{1-\theta}\right)^{1/\nu}(1-\beta)^{1-\frac{1}{\nu}}
\]

\[
\Rightarrow\quad
\omega
=
\chi \eta (\nu-\sigma)
=
\frac{k_m(1-\beta)}{1+k_m(1-\beta)}
\cdot
\frac{1}{\nu\bigl(e^{\bar{i}}-1\bigr)}
\cdot
(\nu-\sigma)
\]

\[
e^{\bar{i}}=\beta^{-1}
\;\Rightarrow\;
e^{\bar{i}}-1=\frac{1-\beta}{\beta}
\]

\[
\Rightarrow\quad
\omega
=
\frac{k_m\beta\left(\nu-\sigma\right)}
{\nu\left[1+k_m(1-\beta)\right]}
\]
\[
\operatorname{sign}(\omega)=\operatorname{sign}\!\left(1-\frac{\sigma}{\nu}\right)
\qquad\Rightarrow\qquad
\omega>0
\iff
\nu>\sigma
\]

\[
\eta=\frac{1}{\nu\bigl(e^{\bar{i}}-1\bigr)}
\]

\[
\text{for small }\bar{i},\qquad
e^{\bar{i}}-1\approx \bar{i}
\;\Rightarrow\;
\eta\approx \frac{1}{\nu\bar{i}}
\;\Rightarrow\;
\nu\approx \frac{1}{\eta\bar{i}}
\]

\[
\text{since }\bar{i}\text{ is small, }\;
\nu=\frac{1}{\eta\bar{i}}\gg \sigma
\;\Rightarrow\;
\omega>0
\]

\[
\left\{
\begin{aligned}
\hat{w}_t-\hat{p}_t &= \sigma\hat{c}_t+\varphi\hat{n}_t+\omega \hat{i}_t,\\
\hat{m}_t-\hat{p}_t &= \hat{y}_t-\eta \hat{i}_t,
\end{aligned}
\right.
\qquad
\hat{c}_t-\hat{x}_t
=
\chi\bigl[\hat{c}_t-(\hat{m}_t-\hat{p}_t)\bigr]
\]

\[
{\color{red}
\hat{i}_t\uparrow,\ \eta>0
\;\Rightarrow\;
\hat{m}_t-\hat{p}_t\downarrow
\;\Rightarrow\;
\hat{x}_t\downarrow
\;\Rightarrow\;
U_{C,t}\downarrow
\;\Rightarrow\;
U_{C\mathbb{M}}>0
}
\]

\[
\textcolor{red}{\text{Money and consumption are complements.}}
\]

\[
\textcolor{red}{\hat{i}_t\uparrow,\ \eta>0
\;\Rightarrow\;
\hat{n}_t\downarrow
\quad\text{given real wage and consumption}}
\]

\[
\textcolor{red}{
\begin{aligned}
&\text{Money and labor move in the same direction: Less liquidity, } U_{C,t}\downarrow,\\
&\text{consumption is less valuable, and the benefit of working is lower.}
\end{aligned}
}
\]


\paragraph{Log-linearize:}
\[
Q_t
=
\beta E_t\!\left[\frac{U_{C,t+1}}{U_{C,t}}\frac{P_t}{P_{t+1}}\right]
=
\beta E_t\!\left[
\left(\frac{C_{t+1}}{C_t}\right)^{-\nu}
\left(\frac{X_{t+1}}{X_t}\right)^{\nu-\sigma}
\frac{P_t}{P_{t+1}}
\right]
\]

\[
Q_t=e^{-i_t},
\qquad
\beta=e^{-\rho},
\qquad
\Pi_{t+1}=\frac{P_{t+1}}{P_t},
\qquad
\pi_{t+1}=\log \Pi_{t+1}=\log P_{t+1}-\log P_t
\]

\[
\Rightarrow\quad
e^{-i_t}
=
e^{-\rho}
E_t\!\left[
e^{
-\nu(\log C_{t+1}-\log C_t)
+(\nu-\sigma)(\log X_{t+1}-\log X_t)
-\pi_{t+1}
}
\right]
\]

\[
\Rightarrow\quad
\left\{
\begin{aligned}
-i_t
&=
-\rho
+
E_t\!\left[
-\nu(\log C_{t+1}-\log C_t)
+(\nu-\sigma)(\log X_{t+1}-\log X_t)
-\pi_{t+1}
\right],\\[4pt]
-\bar{i}
&=
-\rho
+
E_t\!\left[
-\nu(\log \bar{C}-\log \bar{C})
+(\nu-\sigma)(\log \bar{X}-\log \bar{X})
-\bar{\pi}
\right]
\end{aligned}
\right.
\]

\[
\bar{\pi}=0
\]

\[
\Rightarrow\quad
\hat{i}_t
=
i_t-\bar{i}
=
E_t\!\left[
\nu(\hat{c}_{t+1}-\hat{c}_t)
-(\nu-\sigma)(\hat{x}_{t+1}-\hat{x}_t)
+\pi_{t+1}
\right]
\]

\[
\Rightarrow\quad
\hat{i}_t-E_t\pi_{t+1}
=
\nu E_t(\hat{c}_{t+1}-\hat{c}_t)
-(\nu-\sigma)E_t(\hat{x}_{t+1}-\hat{x}_t)
\]

\[
\text{since}\qquad
\hat{c}_t-\hat{x}_t
=
\chi\bigl[\hat{c}_t-(\hat{m}_t-\hat{p}_t)\bigr]
\]

\[
\hat{m}_t-\hat{p}_t
=
\hat{c}_t-\eta \hat{i}_t
\qquad\Rightarrow\qquad
\hat{c}_t-\hat{x}_t
=
\chi\eta \hat{i}_t
\qquad\Rightarrow\qquad
\hat{x}_t
=
\hat{c}_t-\chi\eta \hat{i}_t
\]

\[
\Rightarrow\quad
\hat{i}_t-E_t\pi_{t+1}
=
\nu E_t\Delta \hat{c}_{t+1}
-(\nu-\sigma)E_t\Delta \hat{c}_{t+1}
+
\textcolor{red}{(\nu-\sigma)\chi\eta}\,E_t\Delta \hat{i}_{t+1}
\]

\[
\Rightarrow\quad
\hat{i}_t-E_t\pi_{t+1}
=
\sigma E_t\hat{c}_{t+1}
-\sigma \hat{c}_t
+
\omega E_t\Delta \hat{i}_{t+1}
\]
\[
\Rightarrow\quad
\sigma E_t\Delta \hat{c}_{t+1}
=
\hat{i}_t-E_t\pi_{t+1}
-\omega E_t\Delta \hat{i}_{t+1}
\]

\[
\Rightarrow\quad
\sigma E_t\hat{c}_{t+1}-\sigma\hat{c}_t
=
\hat{i}_t-E_t\pi_{t+1}
-\omega E_t\Delta \hat{i}_{t+1}
\]

\[
\Rightarrow\quad
\hat{c}_t
=
E_t\hat{c}_{t+1}
-\frac{1}{\sigma}
\left[
\hat{i}_t-E_t\pi_{t+1}
-\omega E_t\Delta \hat{i}_{t+1}
\right]
\qquad \cdots\; (*)
\]

\[
\text{where}\qquad
\omega=(\nu-\sigma)\chi\eta,
\qquad
\chi=\frac{k_m(1-\beta)}{1+k_m(1-\beta)},
\qquad
k_m=\frac{\bar{M}/\bar{P}}{\bar{C}},
\qquad
\eta=\frac{1}{\nu(e^{\bar{i}}-1)}.
\]

\[
\textcolor{red}{\text{In }(*)\text{ we apply market clearing }\hat{c}_t=\hat{y}_t}
\]

\[
\textcolor{red}{\text{households smooth consumption}
\;\longrightarrow\;
\text{how responsive current spending is to intertemporal incentives}}
\]

\[
\Rightarrow\quad
\hat{y}_t
=
E_t\hat{y}_{t+1}
-\frac{1}{\sigma}
\left[
\hat{i}_t-E_t\pi_{t+1}
-\omega E_t\Delta \hat{i}_{t+1}
\right]
\]


\[
\textcolor{red}{
\begin{aligned}
&\text{current demand/output} \\
&= \text{expected future output/demand}
-\text{IES}\times \text{expected real interest rate}
-\text{forward guidance}
\end{aligned}
}
\]
\[
\textcolor{red}{
\text{For }\nu>\sigma\Rightarrow\omega>0.
\ \text{When }E_t\Delta\hat{i}_{t+1}\uparrow,\ \hat{c}_t\uparrow.
\ \text{forward guidance}
}
\]

\textcolor{red}{If \(E_t\Delta \hat{i}_{t+1}\uparrow\), then future money holdings become less attractive, future liquidity services decline, and households shift some consumption from the future to today; hence \(\hat{c}_t\uparrow\).}

\[
\frac{W_t}{P_t}
=
(1-\alpha)A_tN_t^{-\alpha},
\qquad
Y_t=A_tN_t^{1-\alpha}
\]

\[
\hat{w}_t-\hat{p}_t=\hat{a}_t-\alpha \hat{n}_t,
\qquad
\hat{y}_t=\hat{a}_t+(1-\alpha)\hat{n}_t
\]

\[
\Rightarrow\quad
\hat{n}_t=\frac{1}{1-\alpha}(\hat{y}_t-\hat{a}_t)
\]

\[
\Rightarrow\quad
\hat{w}_t-\hat{p}_t=\hat{y}_t-\hat{n}_t
\]

\[
\hat{w}_t-\hat{p}_t=\sigma\hat{c}_t+\varphi \hat{n}_t+\omega \hat{i}_t
\]
\[
\Rightarrow\quad
\hat{y}_t-\hat{n}_t
=
\sigma \hat{y}_t+\varphi \hat{n}_t+\omega \hat{i}_t
\]

\[
\Rightarrow\quad
(1-\sigma)\hat{y}_t-(1+\varphi)\hat{n}_t
=
\omega \hat{i}_t
\]

\[
\Rightarrow\quad
\hat{n}_t
=
\frac{1-\sigma}{1+\varphi}\hat{y}_t
-
\frac{\omega}{1+\varphi}\hat{i}_t
\]

\[
\hat{y}_t
=
\hat{a}_t+(1-\alpha)\hat{n}_t
\]

\[
\Rightarrow\quad
\hat{y}_t
=
\hat{a}_t
+
(1-\alpha)
\left[
\frac{1-\sigma}{1+\varphi}\hat{y}_t
-
\frac{\omega}{1+\varphi}\hat{i}_t
\right]
\]

\[
\Rightarrow\quad
\hat{y}_t
=
\frac{1+\varphi}{\alpha(1-\sigma)+\sigma+\varphi}\hat{a}_t
-
\frac{(1-\alpha)\omega}{\alpha(1-\sigma)+\sigma+\varphi}\hat{i}_t
\qquad (1)
\]
\[
\Rightarrow\quad
\hat{y}_t-\hat{n}_t
=
\sigma\hat{y}_t+\varphi \hat{n}_t+\omega \hat{i}_t
\]

\[
\Rightarrow\quad
(1-\sigma)\hat{y}_t-(1+\varphi)\hat{n}_t=\omega \hat{i}_t
\]

\[
\Rightarrow\quad
\hat{n}_t
=
\frac{1-\sigma}{1+\varphi}\hat{y}_t
-
\frac{\omega}{1+\varphi}\hat{i}_t
\]

\[
\hat{y}_t=\hat{a}_t+(1-\alpha)\hat{n}_t
\]

\[
\Rightarrow\quad
\hat{y}_t
=
\hat{a}_t+(1-\alpha)
\left[
\frac{1-\sigma}{1+\varphi}\hat{y}_t
-
\frac{\omega}{1+\varphi}\hat{i}_t
\right]
\]

\[
\Rightarrow\quad
\hat{y}_t
=
\frac{1+\varphi}{\alpha(1-\sigma)+\sigma+\varphi}\hat{a}_t
-
\frac{(1-\alpha)\omega}{\alpha(1-\sigma)+\sigma+\varphi}\hat{i}_t
\qquad (1)
\]

\[
\textcolor{red}{\hat{y}_t:
\text{ interest rate }\to \text{ real activity }\to
\text{transmission coefficient}}
\]

\[
\textcolor{red}{\text{comes from money-in-the-utility non-separability.}}
\]

\[
\text{Money is non-neutral}
\qquad
\textcolor{red}{(\text{interest rate / monetary policy affects real output})}
\]

\[
\text{But the function alone cannot pin down it; we need a monetary policy rule.}
\]

\[
\text{Write the system of equations in terms of }
(\hat{y}_t,\pi_t,\hat{i}_t).
\]

\[
\text{Euler equation:}\qquad
\hat{y}_t
=
E_t\hat{y}_{t+1}
-\frac{1}{\sigma}
\left[
\hat{i}_t-E_t\pi_{t+1}
-\omega E_t\Delta \hat{i}_{t+1}
\right]
\qquad (2)
\]

\[
\textcolor{red}{\text{Also known as IS equation}}
\]

\[
\text{Monetary policy rule:}\qquad
\hat{i}_t=\phi_\pi \pi_t+v_t
\qquad (3)
\]

\[
\text{where}\qquad
v_t=\rho_v v_{t-1}+\varepsilon_t^v
\qquad (4)
\]

\[
\hat{a}_t=\rho_a \hat{a}_{t-1}+\varepsilon_t^a
\qquad (5)
\]

\[
\text{Combine (1)--(5), use undetermined coefficients:}
\]

\[
\pi_t = A_{\pi}\hat{a}_t + B_{\pi}v_t,
\qquad
\hat{i}_t = A_i \hat{a}_t + B_i v_t,
\qquad
\hat{y}_t = A_y \hat{a}_t + B_y v_t.
\]
\[
\Rightarrow\quad
\left\{
\begin{aligned}
{\pi}_t
&=
-\frac{\sigma(1-\rho_a)\psi_{ya}}
{\phi_{\pi}(1+\omega\psi)(1-\theta\rho_a)}\,\hat{a}_t
-\frac{1+(1-\rho_v)\omega\psi}
{\phi_{\pi}(1+\omega\psi)(1-\theta\rho_v)}\,v_t,
\\[8pt]
\hat{i}_t
&=
-\frac{\sigma(1-\rho_a)\psi_{ya}}
{(1+\omega\psi)(1-\theta\rho_a)}\,\hat{a}_t
-\frac{\rho_v}
{\phi_n(1+\omega\psi)(1-\theta\rho_v)}\,v_t,
\\[8pt]
\hat{y}_t
&=
\psi_{ya}\left(
1+
\frac{\sigma(1-\rho_a)\psi_{yi}}
{(1+\omega\psi)(1-\theta\rho_a)}
\right)\hat{a}_t
+
\frac{\rho_v\psi_{yi}}
{\phi_{\pi}(1+\omega\psi)(1-\theta\rho_v)}\,v_t.
\end{aligned}
\right.
\]
\[
\text{where}\qquad
\theta
=
\frac{1+\omega\varphi\phi_{\pi}}
{(1+\omega\varphi)\phi_{\pi}},
\qquad
\psi
=
\frac{\alpha+\varphi}{\sigma(1-\alpha)+\alpha+\varphi}.
\]

\[
\textcolor{red}{\text{Intuition: }}
\psi
\textcolor{red}{\text{ is a summary slope coming from the real-side block, capturing how strongly}}
\]
\[
\textcolor{red}{\text{real activities / labor market tightness maps into wages and output in the log-linear world.}}
\]

\noindent
$\theta$ is a dynamic feedback object mixing the Taylor coefficient $\phi_{\pi}$,
the forward-guidance strength $\omega$, and the real-side slope.

\noindent
When $\omega=0$, $\theta$ collapses to $1/\phi_{\pi}$, which is the usual policy-feedback scale.

\noindent
When $\omega\neq 0$, $\theta$ changes because expected changes in it enter the IS equation.

\paragraph{Impact of the monetary policy}


\[
\text{Conditional on an exogenous MP shock }\varepsilon_t^v\text{ to }v_t:
\qquad
\hat{y}_t=\psi_{ya}\hat{a}_t-\psi_{yi}\hat{i}_t
\qquad\Rightarrow\qquad
\frac{d\hat{y}_t}{d\hat{i}_t}=-\psi_{yi}.
\]

\[
\text{Assume}\qquad
\sigma=\phi=1,\qquad \alpha=\frac{1}{3}
\qquad\Rightarrow\qquad
\psi_{yi}=\frac{\omega}{3}.
\]

\[
\omega
=
\frac{k_m\beta\left(1-\frac{\sigma}{\nu}\right)}
{1+k_m(1-\beta)}
\approx
k_m
\qquad
\text{with }\nu\text{ large and }\beta\approx 1.
\]

\[
\Rightarrow\qquad
\psi_{yi}\approx \frac{k_m}{3}.
\]

\[
\text{Recall}\qquad
k_m=\frac{\bar{M}/\bar{P}}{\bar{C}},
\]
\[
\text{how much liquidity stock households want to hold per unit of spending(cash buffer relative to spending)}
\]

\noindent
Using M1, the money with highest liquidity: \quad
$k_m \approx 0.3 \;\Longrightarrow\; \psi_{yi} \approx 0.1$.

\medskip

\noindent
If $i_t$ increases by 1\% annually (0.25\% quarterly), $y_t$ decreases by 0.025\% quarterly.

\bigskip

\noindent
Using M2, the money with highest liquidity + 2nd highest liquidity: \quad
$k_m \approx 3 \;\Longrightarrow\; \psi_{yi} \approx 1$.

\medskip

\noindent
If $i_t$ increases by 1\% annually, $y_t$ decreases by 1\% annually.

\[
\textcolor{red}{\text{* If money is quantitatively large, MIU channel becomes stronger.}}
\]
\paragraph{Deficiency of MIU model}

\noindent
\[
\frac{d\pi_t}{d\hat{i}_t}
=
\frac{d\pi_t/d\nu_t}{d\hat{i}_t/d\nu_t}
=
\bigl(1+(1-\rho_{\nu})\omega\psi\bigr)\rho_{\nu}^{-1}
>0,
\]
\[
\frac{d\hat{r}_t}{d\hat{i}_t}
=
1-
\frac{dE_t(\pi_{t+1})/d\nu_t}{d\hat{i}_t/d\nu_t}
=
-(1-\rho_{\nu})\omega\psi
<0.
\]

\noindent
A monetary policy shock that raises the nominal interest rate and lowers output also results in:

\noindent
(i) inflation tending to increase;

\noindent
(ii) the real interest rate tending to decrease
\[
\bigl({r}_t={i}_t-E_t\pi_{t+1}\bigr).
\]

\noindent
This contradicts the data: contractionary monetary policy shocks raise real interest rates and lower inflation.

\noindent
In the data, money is neutral in the long run once prices adjust, but non-neutral in the short run. In the model, long-run and short-run effects are the same.

\paragraph{Uniqueness of the nominal system}\mbox{}\\

\noindent
In a flexible-price (classical) economy, the real side pins down the real interest rate $\hat{r}_t$.
The question is: what must monetary policy do to ensure that inflation $\pi_t$
--- and hence the price level $p_t$ --- is uniquely determined?

\noindent
The Fisher relation is
\[
\hat{i}_t = E_t[\pi_{t+1}] + \hat{r}_t,
\qquad\Rightarrow\qquad
E_t[\pi_{t+1}] = \hat{i}_t - \hat{r}_t.
\]

\noindent
Since $\hat{r}_t$ is pinned down by the real economy, if the central bank exogenously sets
$\hat{i}_t$ as some stationary process, or simply sets $\hat{i}_t=0$, then the Fisher equation
pins down expected inflation only, not realized inflation.

\noindent
Therefore, monetary policy must do more than fix the nominal interest rate exogenously:
it must provide a rule that anchors the nominal system and uniquely determines inflation.

\[
\Rightarrow\qquad
E_t[\pi_{t+1}]=-\hat{r}_t
\qquad\Rightarrow\qquad
\pi_{t+1}=-\hat{r}_t+\zeta_{t+1},
\qquad
E_t[\zeta_{t+1}]=0.
\]

\[
\text{Any zero-mean innovation can be added without violating Fisher equation}
\;\Rightarrow\;
\text{indeterminacy.}
\]

\[
\text{With an exogenous }\hat{i}_t,\text{ there are infinitely many inflation paths consistent with equilibrium.}
\]

\[
\Rightarrow\qquad
\text{sunspot driven (non-fundamental factors).}
\]


\[
\text{Since}\qquad
\hat{p}_t=\pi_t+\hat{p}_{t-1},
\qquad
\pi_t\text{ is indetermined}
\;\Rightarrow\;
\hat{p}_t\text{ is also indetermined (not unique).}
\]

\[
\text{Other nominal variables are also indetermined:}
\]

\[
\hat{m}_t=\hat{p}_t+\hat{y}_t-\eta \hat{i}_t
\qquad\Rightarrow\qquad
\hat{m}_t\text{ indetermined}
\qquad
(\text{nominal money})
\]

\[
\hat{w}_t=\hat{w}_t^{r}+\hat{p}_t
\qquad\Rightarrow\qquad
\hat{W}_t\text{ indetermined}
\qquad
(\text{nominal wage})
\]

\noindent
At steady state,
\[
P_t=\bar{P}
\;\Rightarrow\;
\bar{\pi}=0
\;\Rightarrow\;
\bar{i}=\bar{r}=\rho,
\]
so that
\[
\hat{r}_t = r_t - \rho.
\]

\medskip

\noindent
A simple inflation-based interest rate rule (Taylor style) is
\[
\hat{i}_t=\phi_{\pi}\pi_t,
\qquad \phi_{\pi}\ge 0.
\]

\noindent
\[
\hat{i}_t=E_t[\pi_{t+1}] + \hat{r}_t,
\]
we obtain
\[
\phi_{\pi}\pi_t = E_t[\pi_{t+1}] + \hat{r}_t.
\tag{*}
\]

\noindent
This is a forward-looking difference equation in inflation.

\noindent
\textbf{Case 1:}\quad $\phi_{\pi}>1$

\noindent
Iterate forward:
\[
\pi_t=\frac{1}{\phi_{\pi}}E_t[\pi_{t+1}]+\frac{1}{\phi_{\pi}}\hat r_t
\quad\Rightarrow\quad
\pi_t=\sum_{k=0}^{\infty}\phi_{\pi}^{-(k+1)}E_t[\hat r_{t+k}],
\]
where the forward sum converges and yields a unique bounded solution.

\smallskip

\noindent
Bond Euler:
\[
\hat i_t-E_t[\pi_{t+1}]=\sigma E_t\Delta \hat c_{t+1},
\qquad
\hat i_t-E_t[\pi_{t+1}]=\hat r_t,
\]
so
\[
\hat r_t=\sigma E_t\Delta \hat c_{t+1}.
\]

\smallskip

\noindent
Since \(\Delta \hat c_{t+1}=\Delta \hat y_{t+1}\), and
\[
\hat y_t=\psi_{ya}\hat a_t
\quad\Rightarrow\quad
\Delta \hat y_{t+1}=\psi_{ya}\Delta \hat a_{t+1},
\]
we obtain
\[
\hat r_t=\sigma\psi_{ya}E_t\Delta \hat a_{t+1}.
\]

\smallskip

\noindent
If productivity follows
\[
\hat a_{t+1}=\rho_a\hat a_t+\varepsilon^a_{t+1},
\]
then
\[
E_t[\hat a_{t+1}]=\rho_a\hat a_t,
\qquad
E_t[\Delta \hat a_{t+1}]
=
E_t[\hat a_{t+1}-\hat a_t]
=
(\rho_a-1)\hat a_t.
\]

\noindent
Hence,
\[
\hat r_t=\sigma\psi_{ya}(\rho_a-1)\hat a_t.
\]
\smallskip

\noindent
Guess a linear solution:
\[
\pi_t=K\hat a_t,
\qquad
E_t\hat a_{t+1}=\rho_a\hat a_t.
\]
Substituting into
\[
\phi_{\pi}\pi_t=E_t\pi_{t+1}+\hat r_t
\]
gives
\[
\phi_{\pi}K\hat a_t=K\rho_a\hat a_t+\hat r_t
=K\rho_a\hat a_t-\sigma\psi_{ya}(1-\rho_a)\hat a_t,
\]
so
\[
\phi_{\pi}K-K\rho_a=-\sigma\psi_{ya}(1-\rho_a),
\]
and therefore
\[
K=-\frac{\sigma\psi_{ya}(1-\rho_a)}{\phi_{\pi}-\rho_a}.
\]
Thus,
\[
\pi_t=-\frac{\sigma\psi_{ya}(1-\rho_a)}{\phi_{\pi}-\rho_a}\hat a_t.
\]

\smallskip

\noindent
The larger the $\phi_{\pi}$, the smaller the variation of inflation caused by $\hat a_t$.
Central bank raises it more when inflation rises.
Any shock that pushes up $\pi_t$ implies a larger rise in $i_t$; through the Fisher equation,
expectations are forced to be consistent only with the saddle path.

\bigskip

\paragraph{Case 2: $\phi_{\pi}\le 1$}

\noindent
Forward iteration breaks. Since $\phi_{\pi}^{-1}\ge 1$, boundedness no longer pins down a unique $\pi_t$.

\smallskip

\noindent
To satisfy the Fisher equation,
\[
\phi_{\pi}\pi_t=E_t[\pi_{t+1}]+\hat r_t.
\]
The stable solution is
\[
\pi_{t+1}=\phi_{\pi}\pi_t-\hat r_t+\xi_{t+1},
\qquad
E_t[\xi_{t+1}]=0,
\]
with sunspot shocks. Inflation and the price level are indetermined.

\smallskip

\noindent
If the central bank responds less than one-for-one to inflation(weak response), expected inflation can move around
without being forcefully pushed back by policy. This leaves room for self-fulfilling beliefs about inflation
(sunspot shocks), just as in the exogenous $i_t$ case.

\smallskip

\noindent
Determinacy requires the Taylor principle: raise the nominal rate more than one-for-one with inflation,
that is, $\phi_{\pi}>1$.

\smallskip

\noindent
The system is
\[
E_t\pi_{t+1}=\phi_{\pi}\pi_t-\hat r_t.
\]
Inflation $\pi_t$ is a jump variable. By the Blanchard--Kahn condition, to pin down a unique equilibrium
we need
\[
\#\text{ unstable roots}=\#\text{ jump variables}.
\]
From
\[
E_t\pi_{t+1}=\phi_{\pi}\pi_t,
\]
the eigenvalue is $\phi_{\pi}$. If $\phi_{\pi}>1$, there is an unstable root, so the BK condition is satisfied.
If $\phi_{\pi}\le 1$, there is no unstable root to discipline the jump variable, and we get many bounded paths:
sunspot indeterminacy.

\bigskip

\paragraph{Optimal Monetary Policy}

\noindent
Assume a social planner takes the role of the central banker.

\smallskip

\noindent
Social planner's problem:
\[
\max \; U\!\left(C_t,\frac{M_t}{P_t},N_t\right)
\qquad
\text{s.t.}
\qquad
C_t=A_tN_t^{1-\alpha}.
\]

\smallskip

\noindent
The Lagrangian is
\[
\mathcal{L}
=
U\!\left(C_t,\frac{M_t}{P_t},N_t\right)
+\lambda_t\bigl(A_tN_t^{1-\alpha}-C_t\bigr).
\]

\smallskip

\noindent
F.O.C.:
\[
-\frac{U_{N,t}}{U_{C,t}}
=
(1-\alpha)A_tN_t^{-\alpha}
=
W_t
=
MPL_t
=
MRS_{N,C},
\]
and
\[
U_{\mathbb{M},t}=0.
\]

\smallskip

\noindent
In competitive equilibrium,
\[
\frac{U_{\mathbb{M},t}}{U_{C,t}}=1-e^{-i_t}\approx i_t.
\]
Social and individual optima coincide iff $i_t=0$ for all $t$, which is the Friedman rule.
\smallskip

\noindent
In the long run, with the Fisher equation,
\[
i=\pi+r=0
\quad\Rightarrow\quad
\pi=-r=-\rho<0,
\]
which is a moderate deflation.

\smallskip

\noindent
In the short run, $i_t=0$ does not guarantee determinacy. Instead,
\[
E_t\pi_{t+1}=-r_t,
\]
or equivalently,
\[
\pi_{t+1}=-r_t+\xi_{t+1},
\qquad
E_t[\xi_{t+1}]=0,
\]
so any sunspot shock can be added.

\smallskip

\noindent
Even though money balances and other nominal terms are indetermined, real variables are pinned down by
\[
MRS=MPL.
\]
Thus indeterminacy does not have welfare consequences.

\smallskip

\noindent
To avoid indeterminacy, let
\[
i_t=\phi\,(r_{t-1}+\pi_t),
\qquad \phi>1.
\]
Then
\[
i_t=E_t\pi_{t+1}+r_t
\]
\[
E_t[i_{t+s}] = \phi^s i_t.
\]
Since \(\phi>1\), as \(s\to\infty\), \(i_t\) has to approach \(0\).
\[
\Rightarrow\quad \pi_t=-r_{t-1}\text{, so inflation is fully determined.}
\]
Any policy setting that guarantees that current inflation moves inversely, and one-for-one, with \(r_{t-1}\) will imply \(i_t=0\) and an efficient amount of real balances.

\bigskip

\paragraph{Supplementary: IS equation under the separable case in the MIU model}

\noindent
\[
U\!\left(C_t,\mathbb{M}_t,N_t\right)
=
\frac{C_t^{1-\sigma}}{1-\sigma}
+
\frac{(\mathbb{M}_t)^{1-\nu}}{1-\nu}
-
\frac{N_t^{1+\varphi}}{1+\varphi}.
\]
\smallskip

\noindent
From the bond Euler F.O.C.,
\[
Q_t=\beta E_t\!\left[\frac{U_{C,t+1}}{U_{C,t}}\frac{P_t}{P_{t+1}}\right].
\]
Since
\[
U_{C,t}=C_t^{-\sigma},
\]
we have
\[
Q_t=\beta E_t\!\left[\left(\frac{C_{t+1}}{C_t}\right)^{-\sigma}\frac{P_t}{P_{t+1}}\right].
\]

\smallskip

\noindent
Define
\[
i_t=-\log Q_t,
\qquad
\rho=-\log\beta,
\qquad
\pi_{t+1}=\log P_{t+1}-\log P_t.
\]
Then
\[
-i_t
=
-\rho
+
E_t\!\left[-\sigma(\log C_{t+1}-\log C_t)-\pi_{t+1}\right].
\]

\smallskip

\noindent
At steady state,
\[
-\bar i
=
-\rho
+
E_t\!\left[-\sigma(\log\bar C-\log\bar C)-\bar\pi\right],
\]
with
\[
\bar\pi=0.
\]
Hence
\[
\bar i=\rho.
\]

\smallskip

\noindent
Subtracting the steady state and using hats,
\[
\sigma E_t(\hat c_{t+1}-\hat c_t)=\hat i_t-E_t\pi_{t+1}.
\]
Thus
\[
\hat c_t
=
E_t\hat c_{t+1}
-
\frac{1}{\sigma}\bigl[\hat i_t-E_t\pi_{t+1}\bigr].
\]
Since
\[
\hat y_t=\hat c_t,
\]
we get
\[
\hat y_t
=
E_t\hat y_{t+1}
-
\frac{1}{\sigma}\bigl[\hat i_t-E_t\pi_{t+1}\bigr].
\]
Compare to the classical monetary model: the same.

Compare to the inseparable case: \(\nu=\sigma\), and the term \(-\omega E_t\Delta \hat{i}_{t+1}\) disappears?
\smallskip


\bigskip
\noindent\textbf{Supplement: Interpret}
\[
\hat{y}_t
=
E_t\hat{y}_{t+1}
-\frac{1}{\sigma}\left[\hat{i}_t-E_t\pi_{t+1}-\omega E_t\Delta\hat{i}_{t+1}\right].
\]
\noindent
1. The term
\[
i_t-E_t\pi_{t+1}
\]
is the opportunity cost of consuming today instead of tomorrow.

\smallskip

\noindent
Suppose you have 1 RMB.

\smallskip

\noindent
If you buy goods today, you get
\[
\frac{1}{P_t}
\]
units of goods.

\smallskip

\noindent
If you save to tomorrow and then buy goods, you get
\[
\frac{1+i_t}{P_{t+1}}
\]
units of goods. Therefore,
\[
E_t\!\left[\frac{(1+i_t)/P_{t+1}}{1/P_t}\right]
=
E_t\!\left[\frac{1+i_t}{P_{t+1}/P_t}\right]
=
\frac{1+i_t}{1+E_t\pi_{t+1}}
\approx
i_t-E_t\pi_{t+1}.
\]

\smallskip

\noindent
So if
\[
\hat i_t-E_t\pi_{t+1}\uparrow,
\]
then
\[
\hat y_t\downarrow,
\]
because the opportunity cost of consuming today is higher.

\smallskip

\noindent
2. The term
\[
\omega E_t\hat i_{t+1}
\]
is the forward-guidance term.

\smallskip

\noindent
If tomorrow's nominal rate $i_{t+1}$ is expected to be high, then holding money tomorrow is expensive.
Households therefore want to hold fewer real balances at $t+1$:
\[
\frac{M_{t+1}}{P_{t+1}}\downarrow.
\]

\smallskip

\noindent
With $\omega>0$, there is within-period substitution between $C_t$ and $\frac{M_{t+1}}{P_{t+1}}$,
in addition to the usual intertemporal substitution of $X_t$.

\smallskip

\noindent
Holding more real balances,
\[
\frac{M_{t+1}}{P_{t+1}}\uparrow,
\]
increases the marginal utility of consumption, $MU_{C,t+1}$, so future consumption becomes more valuable.
Conversely,
\[
\frac{M_{t+1}}{P_{t+1}}\downarrow
\]
makes future consumption less valuable and induces households to consume more today.

\smallskip

\noindent
In particular,
\[
U_{C,t}=(1-\theta)X_t^{\nu-\sigma}C_t^{-\nu},
\qquad
X_t \propto \mathbb{M}_t
\quad\Rightarrow\quad
U_{C\mathbb{M}}>0.
\]

\smallskip

\noindent
Consumption is more enjoyable with more liquidity.

\bigskip

\paragraph{Supplement for computation:Log-linearized bond Euler equation}
\noindent
Define
\[
i_t=-\log Q_t \;\Rightarrow\; e^{-i_t}=Q_t,
\qquad
\rho=-\log\beta \;\Rightarrow\; e^{-\rho}=\beta,
\qquad
\pi_{t+1}=\log P_{t+1}-\log P_t,
\]
and
\[
r_t=i_t-E_t[\pi_{t+1}].
\]

\noindent
Define hats as level deviations:
\[
\hat i_t=i_t-\bar i,
\qquad
\hat\pi_{t+1}=\pi_{t+1}-\bar\pi,
\qquad
\bar\pi=0 \;\Rightarrow\; \hat\pi_{t+1}=\pi_{t+1}.
\]
Also,
\[
r_t-\bar r=i_t-\bar i-E_t[\pi_{t+1}-\bar\pi]
\;\Rightarrow\;
\hat r_t=\hat i_t-E_t[\hat\pi_{t+1}].
\]

\noindent
From the bond Euler equation,
\[
Q_t=\beta E_t\!\left[\left(\frac{C_{t+1}}{C_t}\right)^{-\sigma}\frac{P_t}{P_{t+1}}\right]
=\beta E_t\!\left[e^{-\sigma(\log C_{t+1}-\log C_t)-\pi_{t+1}}\right],
\]
so
\[
e^{-i_t}=e^{-\rho}E_t\!\left[e^{-\sigma(\log C_{t+1}-\log C_t)-\pi_{t+1}}\right].
\]

\noindent
Let
\[
\hat c_{t+1}=\log C_{t+1}-\log\bar C,
\qquad
\hat c_t=\log C_t-\log\bar C
\;\Rightarrow\;
\log C_{t+1}-\log C_t=\hat c_{t+1}-\hat c_t.
\]
Then
\[
e^{-i_t}=e^{-\rho}E_t\!\left[e^{-\sigma(\hat c_{t+1}-\hat c_t)-\pi_{t+1}}\right].
\]

\noindent
Taking logs,
\[
-i_t=-\rho+E_t\!\left[-\sigma(\hat c_{t+1}-\hat c_t)-\pi_{t+1}\right].
\]
At steady state,
\[
-\bar i=-\rho+E_t[0-0] \;\Rightarrow\; \bar i=\rho.
\]
Hence
\[
-(i_t-\bar i)=E_t\!\left[-\sigma(\hat c_{t+1}-\hat c_t)\right]-E_t[\pi_{t+1}-\bar\pi],
\]
or equivalently,
\[
\hat i_t=\sigma E_t[\Delta \hat c_{t+1}]+E_t[\hat\pi_{t+1}],
\qquad
\Delta \hat c_{t+1}=\hat c_{t+1}-\hat c_t.
\]
Therefore,
\[
\sigma E_t[\Delta \hat c_{t+1}]
=\hat i_t-E_t[\hat\pi_{t+1}]
=\hat r_t.
\]
Using \(\hat c_t=\hat y_t\), we obtain
\[
\sigma E_t[\Delta \hat y_{t+1}]=\hat r_t.
\]

\noindent
Since
\[
E_t[\Delta \hat y_{t+1}]=E_t\hat y_{t+1}-\hat y_t,
\qquad
\hat y_t=\psi_{ya}\hat a_t
\;\Rightarrow\;
E_t[\Delta \hat y_{t+1}]=\psi_{ya}E_t[\Delta \hat a_{t+1}],
\]
it follows that
\[
\hat r_t=\sigma\psi_{ya}E_t[\Delta \hat a_{t+1}].
\]

\noindent
Let
\[
\log A_{t+1}=(1-\psi)\log\bar A+\varphi\log A_t+\varepsilon_{t+1},
\qquad
\bar A=1.
\]
Then
\[
\log\bar A=\psi\log\bar A,
\qquad
\hat a_{t+1}=\psi\hat a_t+\varepsilon_{t+1}.
\]
Taking expectations,
\[
E_t[\hat a_{t+1}]
=\psi\hat a_t+E_t[\varepsilon_{t+1}],
\qquad
\varepsilon_t\sim \text{i.i.d. }N(0,\sigma^2)
\;\Rightarrow\;
E_t[\varepsilon_{t+1}]=0,
\]
so
\[
E_t[\hat a_{t+1}]=\psi\hat a_t,
\qquad
E_t[\hat a_{t+1}-\hat a_t]=(\psi-1)\hat a_t.
\]
Let
\[
\Delta \hat a_{t+1}=\hat a_{t+1}-\hat a_t
\;\Rightarrow\;
E_t[\Delta \hat a_{t+1}]=(\psi-1)\hat a_t.
\]

\noindent
Plugging this into \(\hat r_t=\sigma\psi_{ya}E_t[\Delta \hat a_{t+1}]\), we get
\[
\hat r_t
=\sigma\psi_{ya}E_t[\Delta \hat a_{t+1}]
=\sigma\psi_{ya}(\psi-1)\hat a_t.
\]
Since \(\psi\in(0,1)\) and \(\psi_{ya}>0\), the coefficient on \(\hat a_t\) is negative.

\ThisNoteEnd
